% LSPV-1 REPORT::HEVALLB1::2026-10-12T05:09:35Z
% !TEX program = xelatex
\documentclass[11pt]{article}
\usepackage[margin=0.84in]{geometry}
\usepackage{amsmath,amssymb,amsthm,booktabs,array,enumitem,xurl,hyperref}
\setlength{\emergencystretch}{3em}
\hypersetup{colorlinks=true,linkcolor=blue,urlcolor=blue}
\newtheorem{definition}{Definition}
\newtheorem{theorem}{Theorem}
\newtheorem{corollary}[theorem]{Corollary}
\newcommand{\SAT}{\mathsf{SAT}}
\newcommand{\Hsigned}{\mathcal H^{\pm}_{2,2}}
\newcommand{\EXACT}{\textbf{PROVED/EXACT}}
\newcommand{\KNOWN}{\textbf{KNOWN THEOREM}}
\newcommand{\COMP}{\textbf{COMPUTATIONAL}}
\newcommand{\COND}{\textbf{CONDITIONAL}}
\newcommand{\OPEN}{\textbf{OPEN}}
\newcommand{\RETIRED}{\textbf{FALSIFIED/RETIRED}}
\title{HEVALLB1: Uniform evaluation of the full partial-assignment\\
compatibility function is exactly SAT-in-P strength}
\author{LSPV-1 --- Universal Core Lower Bound}
\date{12 October 2026 (UTC)}

\begin{document}
\maketitle

\begin{center}
\fbox{\parbox{0.94\linewidth}{
\textbf{Executive status.} P versus NP remains open.  This run fixes a
canonical semantic compatibility object and proves its exact complexity
boundary.  For a CNF $F$ and a partial assignment $\alpha$, let
$H(F,\alpha)=1$ iff $\alpha$ extends to a satisfying assignment.  A single
uniform polynomial-time evaluator for all pairs $(F,\alpha)$ exists if and
only if $\SAT\in\mathsf P$; the same equivalence holds when $F$ is restricted
to any polynomial-time many-one NP-hard clean core, including the corrected
signed Monotone $3$SAT-$(2,2)$ family.  Under $\SAT\in\mathsf P$, $H_F$ has a
linear-size program descriptor and polynomial-size uniformly constructible
circuits, but their exponent inherits the unknown SAT exponent.  Therefore
the $3^n$ semantic domain or explicit truth-table size cannot force
materialization, memory, communication, or runtime for an arbitrary one-bit
SAT execution.  The evaluator is an external new execution, not a
SAME\_REALIZATION exposure theorem.}}
\end{center}

\section{Claim card and exact quantifiers}

\begin{description}[leftmargin=3.0cm,style=nextline]
\item[Claim ID.] \texttt{PARTIAL-COMPAT-EVAL-EQUIV1}.
\item[Status before run.] \OPEN; fixed-split compatibility was already known
not to transfer automatically to one-bit SAT computation.
\item[Desired conclusion.] Classify the complexity of evaluating the complete
partial-assignment compatibility function without enumerating it.
\item[Exact prefix.] $\exists E\,\exists c\,\forall F\,\forall\alpha$ for the
upper-bound side; the clean-core reverse direction also quantifies over one
fixed polynomial many-one reduction into the core.
\item[Model.] Uniform deterministic multitape Turing machine on a complete
self-delimiting encoding.  Malformed inputs may be rejected.  All restriction,
encoding and output work is charged.
\item[Input measure.] $L=|\operatorname{enc}(F,\alpha)|$; for the canonical
encoding $L=N+O(n)$ and $n\le N+O(1)$.
\item[SAME\_REALIZATION.] \texttt{EVAL(F,alpha)} is a new two-input execution.
It is not identified with the original one-input run $A(F)$.
\item[Weakest useful theorem.] An iff classification, plus a precise account
of what compact descriptors and circuits do not transfer back to $A(F)$.
\item[Falsifier.] A family on which $H$ cannot be evaluated by restriction
followed by a SAT decider, or a proof that one-bit SAT forces the whole table
to be output.  Neither occurs.
\end{description}

\section{Canonical semantic compatibility}

\begin{definition}[Partial-assignment compatibility]
Let $F$ be a CNF on indexed variables $x_1,\ldots,x_n$ and let
$\alpha\in\{0,1,\star\}^n$.  Write $F\upharpoonright\alpha$ for the CNF
obtained by deleting clauses satisfied by $\alpha$ and deleting falsified
literals from the remaining clauses.  Define
\[
 H(F,\alpha)=H_F(\alpha)=1
 \quad\Longleftrightarrow\quad
 \exists\beta\in\{0,1\}^n\ [\beta\succeq\alpha\ \wedge\ F(\beta)=1].
 \tag{1}
\]
Equivalently, $H_F(\alpha)=\SAT(F\upharpoonright\alpha)$.  The table of
$H_F$ has $3^n$ entries, but (1) does not prescribe a representation.
\end{definition}

The empty partial assignment is denoted $\star^n$.  It gives the exact
one-bit boundary
\[
 H_F(\star^n)=\SAT(F).                                      \tag{2}
\]

\begin{theorem}[Uniform evaluator equivalence]
\label{thm:eval}
For canonically encoded CNFs, the following are equivalent.
\begin{enumerate}[label=(\roman*)]
\item $\SAT\in\mathsf P$.
\item There are a uniform deterministic machine $E$ and a constant $c$ such
that for every valid $F$ and every $\alpha\in\{0,1,\star\}^{n(F)}$,
\[
 E(\operatorname{enc}(F,\alpha))=H_F(\alpha),\qquad
 T_E(F,\alpha)\le |\operatorname{enc}(F,\alpha)|^c.       \tag{3}
\]
\end{enumerate}
Moreover, let $K$ be any canonically encoded SAT promise family admitting a
polynomial-time many-one reduction $r:\SAT\le_m^p \SAT|_K$.  Then (i) is also
equivalent to (3) restricted to $F\in K$.
\end{theorem}

\begin{proof}
Assume (i), and fix a polynomial-time SAT decider $A$.  On $(F,\alpha)$,
$E$ validates the indices and the trit encoding, constructs
$F\upharpoonright\alpha$ in one scan, and runs $A$ on that residual.  The
residual has length $O(L)$, so the total time is polynomial in $L$.  Equation
(1) proves correctness.  This charges validation and restriction rather than
treating them as free.

Conversely, assume (ii).  Given $F$, form $\alpha=\star^{n(F)}$ in linear
time and return $E(F,\alpha)$.  Equation (2) gives a polynomial SAT decider.

For the core-restricted form, the forward implication is a restriction of
the preceding evaluator.  In the reverse direction, map an arbitrary SAT
instance $G$ to $F=r(G)\in K$, construct $\star^{n(F)}$, and invoke the core
evaluator.  The reduction, empty assignment and evaluator all have
polynomial cost.  No promise recognition outside the image of $r$ is needed.
\end{proof}

\begin{corollary}[Compact evaluator and circuit representation]
\label{cor:compact}
Suppose a fixed SAT decider $A$ runs in time $p(N)$.  For every $F$, the pair
$\langle A,F\rangle$ is a descriptor of length $N+O_A(1)$ for $H_F$: on input
$\alpha$, construct $F\upharpoonright\alpha$ and run $A$.  Standard
time-to-circuit simulation, with $F$ hardwired, yields a uniformly
constructible Boolean circuit for $H_F$ of size polynomial in
$N+p(O(N))$.  In particular, if $p(N)=N^K$, the size is $N^{O(K)}$.
\end{corollary}

This is \COND{} on $\SAT\in\mathsf P$, not a fixed-exponent contraction.
The exponent depends on the unknown $K$.  A bare assertion that each $H_F$
has a small circuit without a uniform constructor is weaker still: the advice
is indexed by the individual formula $F$, not merely by its length, and does
not by itself place SAT in a standard nonuniform class.

\section{What the theorem falsifies---and what it does not}

The implication
\[
 |\operatorname{dom}H_F|=3^n
 \Longrightarrow
 \text{every SAT execution materializes }\Omega(3^n)\text{ bits}           \tag{4}
\]
is false.  A function can have an exponential domain and a linear descriptor;
under $\SAT\in\mathsf P$, Corollary~\ref{cor:compact} gives exactly such a
descriptor for every $H_F$.  Independently of $P$ versus $NP$, the one-bit
contract asks only for (2), not all table entries.  Thus
\texttt{EXPLICIT-COMPAT-TABLE-NECESSITY1} is \RETIRED.

No polynomial-size OBDD or branching program is claimed for every $H_F$.
Knowledge-compilation languages have representation-specific succinctness
and query tradeoffs.  An OBDD, d-DNNF, communication, proof-system, or
pebbling lower bound can therefore remain correct in its declared model while
having no automatic consequence for unrestricted one-bit computation.

The construction $A\mapsto E_A$, where
\[
 E_A(F,\alpha)=A(F\upharpoonright\alpha),                 \tag{5}
\]
is an exact black-box external extraction theorem.  It does not show that the
trace of $A(F)$ contains, exposes, stores or pays for $H_F$.  Equation (5)
runs $A$ on a different input and may do so separately for every queried
$\alpha$.  Acquisition, repeated calls and output are not free.  Therefore
the arbitrary-decider SAME\_REALIZATION bridge remains \OPEN.

\section{Primary-literature and barrier preflight}

The first obstacle was whether restriction evaluation is a genuine uniform
decision problem or an implicit compilation requirement.  Cook's original
SAT completeness theorem supplies the standard reduction/encoding basis
\cite{Cook71}; Pippenger--Fischer supplies the standard machine/circuit
simulation perspective used only for Corollary~\ref{cor:compact}
\cite{PF79}.  These are \textbf{direct theorem transfers}: their objects are
ordinary uniform computations and Boolean circuits.

The second obstacle was representation dependence.  Darwiche--Marquis
classifies knowledge-compilation targets by succinctness and supported
queries \cite{DM02}.  This is an exact source theorem about compilation
languages, but using an OBDD or d-DNNF lower bound against unrestricted SAT
would be an \textbf{analogy requiring a new transfer theorem}.  FLMV's
time--space exclusion is likewise exact only in its stated random-access
machine model \cite{FLMV05}; the compiler landing the same realization there
is absent.

No literature result found supplies the missing all-decider compiler.
The evaluator equivalence relativizes, so Baker--Gill--Solovay remains a
direct warning \cite{BGS75}; it also survives algebraic-oracle enrichment and
therefore does not evade algebrization \cite{AW09}.

\section{Comparison with the actual manuscript entry points}

Actual \path{main4.tex} is the ARC achievable-region/source-pass entry point.
It records partial promise regions, replay, exact-size split acquisition and
successor-coverage failures.  Actual \path{main5.tex} is the WRAPPER entry
point with the corresponding construction ledger.  Both explicitly preserve
all-input acquisition and SAME\_REALIZATION as open.  Neither defines (1) as
an object that every solver must materialize.  HEVALLB1 does not rename or
strengthen either manuscript: it only proves that an externally evaluable
semantic object remains representation- and realization-distinct.

\section{Bounded reproducible audit}

The deterministic script
\path{experiments/review/semantic_compatibility_evaluator_audit.py} exhausts
all partial assignments for ten formulas: trivial SAT and UNSAT, 2SAT, Horn,
affine XOR, an inconsistent Tseitin cycle, bounded-treewidth path, independent
XOR products, signed 3CNF and unused-variable padding.  For each partial
assignment it compares direct extension search with explicit restriction and
residual SAT, verifies downward closure of incompatibility, and checks
DIRECT, COPY, clause order, literal order, tautological padding and INTERP
variable-permutation wrappers.

The run completed in $0.209065$ seconds: $4{,}836$ partial assignments,
$29{,}016$ wrapper checks and $39{,}998$ assertions, with zero failures.
Every full table used $3^n$ evaluator calls while the SAT answer used the
single empty-partial call.  The raw result is
\path{results/raw/review/2026-10-12_semantic_compatibility_evaluator.json}.
SHA-256 values are
\path{570f55e18bc3995ffbf27e73dae793d2e189783eaecb512079fe340e5abdacca}
for the script and
\path{7936a592a3acec3f5090592312fa8113d55782e929fc96f0bd967d97cf42643d}
for the JSON.  This evidence is \COMP{} and supports no asymptotic inference.

\section{Hostile controls and representation audit}

\begin{center}\small
\begin{tabular}{@{}p{0.25\linewidth}p{0.67\linewidth}@{}}
\toprule
Control & Result\\ \midrule
Horn / 2SAT / affine / Tseitin / bounded treewidth & The same $H$ definition
is valid, but family-specific polynomial algorithms evaluate it without a
generic compatibility table.  A large table would be an optional interface.\\
Independent products & $H_{F\wedge G}(\alpha,\beta)=H_F(\alpha)H_G(\beta)$
on disjoint variables; product-domain growth does not force product-table
materialization.\\
Padding / redundant clauses & COPY, tautologies and unused variables can
inflate encodings or domains without changing the decision bit.\\
Alternate encodings / INTERP & Variable permutation conjugates the evaluator;
it does not create a semantic cost lower bound.\\
DEAD / DELAY / RECOMP & DEAD tables may be skipped, DELAY only inflates a
chosen trace, and RECOMP pays for repeated evaluations.  None is unavoidable.\\
One-bit wrapper & It requests only $H_F(\star^n)$.\\
Compact circuits & Conditional exact upper bound, with exponent inherited
from the assumed SAT decider.\\
Branching programs / OBDDs & No universal compactness claim and no transfer
from a restricted representation lower bound.\\
\bottomrule
\end{tabular}
\end{center}

\section{Barrier and post-quantum ledgers}

\paragraph{Relativization.}
The proof is restriction plus one call to the assumed decision procedure and
the empty-partial reverse call.  The same syntax works relative to an oracle
for the corresponding restriction-closed language.  Hence HEVALLB1 is a guard,
not a nonrelativizing separation step.

\paragraph{Algebrization.}
Low-degree extension access can only make some residual queries easier; no
proof line blocks it.  There is no nonalgebrizing hinge.

\paragraph{Natural Proofs.}
No large constructive property useful against general circuits is asserted.
Trying to turn truth-table or representation rank into such a property would
need a separate Razborov--Rudich audit; HEVALLB1 neither supplies nor evades it.

\paragraph{Unknown exponents and model matching.}
If $A$ takes $N^K$ time, (5) and its circuits retain $N^{O(K)}$ cost.  They do
not yield the solver-exponent-independent contraction required to enter an
FLMV-forbidden region.  The source, residual inputs and repeated evaluator
runs are different realizations from $A(F)$.

\paragraph{Post-quantum techniques.}
Schur squares, quadratic-relation ranks, Hilbert functions and syzygies are
\textbf{FAILED TRANSFERS} here: they concern chosen representations and do
not change Theorem~\ref{thm:eval}.  Collision entropy and special-soundness
forks are \textbf{ANALOGY/CANDIDATE} only; a collision-to-forbidden-witness
theorem is absent.  The LWE-style black-box lesson has one exact but negative
transfer: (5) externally extracts coordinate access to $H$ from a SAT solver,
but produces a new execution and inherits its exponent.  Hash/random-oracle
methods remain inadmissible for an unconditional separation.  Compact
algebraic structure must also be tested as an algorithmic shortcut, as the
affine controls demonstrate.

\section{CORPUS REVIEWED}

The GitHub tree at starting head
\texttt{d437c27d1497b2a80bd87036018770c04297191b} was surveyed.  The run read
root \path{main.tex}, \path{AGENTS.md}, \path{CODEX.md}, the complete
pnp-research skill and mode/protocol references, every current task under
\path{tasks/}, actual \path{main4.tex} and \path{main5.tex}, historical MAIN
SYNTHESIS-0022, newest RESOURCEAXIS1, TESTS-0021, APIFLAT1 and COREEQUIV1,
and materially relevant earlier compatibility/communication and cost-certificate
reports.  The HEC, VSC-HEC, G-PPC-HEC, CRS-PPC, CCSPC and RCCN causal
explorations, governance tree, and cited experimental artifacts were
inspected.  Repository search covered \texttt{H\_F}, semantic compatibility,
truth-table materialization, compact representations and prior retired routes.

The predecessor synthesis snapshot
\path{research/main/2026-10-12_0448_MAIN_SYNTHESIS-0022.tex} was read back
from GitHub.  It is substantive and contains the expected MAIN header and
final verification sentinel; no recovery request was needed.

\section{DUPLICATION / CONTRADICTION AUDIT}

MODELBRIDGE-0006 already defines fixed-split
$H_F^{X\mid Y}(a,b)=\exists z\,F(a,b,z)$ and retires the transfer from its
communication complexity to a fixed one-bit SAT run.  CC-0005 already says
truth-table acquisition need not enumerate the table.  Those results are
adopted, not renamed.  HEVALLB1 adds only the exact all-partial-assignment
uniform evaluator iff, the NP-hard-core restriction, and the conditional
descriptor/circuit corollary with unknown-exponent accounting.

There is no contradiction with APIFLAT1: named APIs can be erased while the
external evaluator (5) remains definable.  There is no contradiction with
COREEQUIV1 or RESOURCEAXIS1: an unconditional superpolynomial evaluator lower
bound on an NP-hard core is another presentation of the separation, while a
conditional polynomial evaluator carries the unknown exponent.

\section{WHAT CHANGED RELATIVE TO ROOT MAIN}

Root MAIN says that no theorem forces compatibility extraction.  HEVALLB1
sharpens this with a canonical $H(F,\alpha)$ and an exact boundary:
uniform polynomial evaluation is equivalent to $\SAT\in\mathsf P$, and
explicit table size is not an intermediate resource lower bound.  The
arbitrary-decider compiler, fixed-exponent contraction and model-matched
forbidden endpoint remain \OPEN.

\section{ADOPTED / REJECTED / RETIRED PRIOR COMPONENTS}

\textbf{Adopted:} MODELBRIDGE-0006's fixed-function/one-run distinction;
CC-0005's semantic-object/machine-state guard; TESTS-0021's zero-interface
one-bit decider; APIFLAT1's API-erasure theorem; COREEQUIV1's separation
boundary; and actual main4/main5 only in their stated partial construction
scope.

\textbf{Rejected:} charging $3^n$ table entries to one evaluation; treating a
small per-formula nonuniform circuit as a uniform SAT algorithm; inferring
ordinary runtime from OBDD, communication, proof or pebbling size; and hiding
the SAT exponent inside ``polynomial.''

\textbf{FALSIFIED/RETIRED:}
\path{EXPLICIT-COMPAT-TABLE-NECESSITY1} and
\path{COMPAT-EVALUATOR-AS-FIXED-CONTRACTION1}.

\section{STALE OR OVERSTATED CLAIMS FOUND}

Any current or future wording that treats an exponential compatibility domain
as an exponential machine obligation is overstated unless it names an output
or interface that requires all entries.  Any claim that ``$P=NP$ makes $H$
compact'' must state the representation: a uniform evaluator and circuits are
compact with solver-dependent exponent; compact OBDDs or uniformly acquired
canonical representations do not follow.  No verified result in the current
root MAIN needs promotion beyond its present status.

\section{Strict status and handoff}

\begin{center}\scriptsize
\begin{tabular}{@{}p{0.47\linewidth}p{0.21\linewidth}p{0.24\linewidth}@{}}
\toprule
Claim & Status & Scope\\ \midrule
\texttt{PARTIAL-COMPAT-EVAL-EQUIV1} & \EXACT & Uniform evaluator, all CNF / NP-hard core\\
Compact program descriptor and circuits & \COND & Assumes $\SAT\in\mathsf P$; exponent inherited\\
Finite ten-family audit & \COMP & $39{,}998$ assertions\\
Explicit-table necessity & \RETIRED & Not forced by one-bit semantics\\
Arbitrary-decider internal exposure & \OPEN & No SAME\_REALIZATION theorem\\
Solver-independent contraction / FLMV landing & \OPEN & Unknown exponent and model transfer\\
$P\ne NP$ & \OPEN & No proof\\
\bottomrule
\end{tabular}
\end{center}

The precise next target is not another compatibility representation.  It is
one externally fixed compiler $C$ with
\[
 \begin{aligned}
 \SAT\in\mathsf P\Longrightarrow{}& C(A)\text{ is the same fully charged}\\
 &\text{ordinary machine in a cited solver-exponent-independent}\\
 &\text{forbidden region}.
 \end{aligned}                                                         \tag{6}
\]
All preprocessing, source access, residual construction, evaluator calls,
replay, reconstruction, output, time and peak space must be charged.  Both
the compiler and the landing in (6) remain \OPEN.

\begin{thebibliography}{9}
\bibitem{Cook71}
S. A. Cook, ``The Complexity of Theorem-Proving Procedures,''
\emph{STOC}, 1971, pp. 151--158.
\url{https://doi.org/10.1145/800157.805047}

\bibitem{PF79}
N. Pippenger and M. J. Fischer, ``Relations Among Complexity Measures,''
\emph{J. ACM} 26(2), 1979, pp. 361--381.
\url{https://doi.org/10.1145/322123.322138}

\bibitem{DM02}
A. Darwiche and P. Marquis, ``A Knowledge Compilation Map,''
\emph{J. Artificial Intelligence Research} 17, 2002, pp. 229--264.
\url{https://www.cs.cmu.edu/afs/cs.cmu.edu/project/jair/pub/volume17/darwiche02a.pdf}

\bibitem{FLMV05}
L. Fortnow, R. Lipton, D. van Melkebeek, and A. Viglas,
``Time-Space Lower Bounds for Satisfiability,'' \emph{J. ACM} 52(6), 2005.
\url{https://doi.org/10.1145/1101821.1101822}

\bibitem{BGS75}
T. Baker, J. Gill, and R. Solovay, ``Relativizations of the $P=?NP$
Question,'' \emph{SIAM J. Comput.} 4(4), 1975, pp. 431--442.
\url{https://doi.org/10.1137/0204037}

\bibitem{AW09}
S. Aaronson and A. Wigderson, ``Algebrization: A New Barrier in Complexity
Theory,'' \emph{ACM TOCT} 1(1), 2009.
\url{https://eccc.weizmann.ac.il/report/2008/034/}
\end{thebibliography}

\noindent\textbf{Final status:} P versus NP remains \OPEN.

% END_OF_VERIFIED_REPORT::LSPV-1::HEVALLB1
\end{document}

